\chapter{Introduction}

\section{Background and Motivation}

Cloud computing has fundamentally transformed how software systems are architected, deployed, and scaled. Modern cloud-native applications increasingly adopt event-driven, microservices-based architectures where loosely coupled components communicate asynchronously through message queues~\citep{Thaur2023EventDrivenKafka, Laigner2024EventChallenges}. This architectural pattern offers significant advantages: it decouples producers from consumers, enables independent scaling of services, provides natural buffering against load spikes, and facilitates fault isolation by preventing cascading failures across synchronous call chains.

Amazon Simple Queue Service (SQS) is one of the most widely adopted managed message queuing services, offering fully managed, serverless infrastructure with high availability, at-least-once delivery guarantees, and seamless integration with AWS Lambda and other compute services~\citep{AWSSQSGuide}. Organizations use SQS for diverse workloads including order processing, log aggregation, event routing, and batch job coordination. The service abstracts the operational complexity of running distributed message brokers, allowing developers to focus on application logic rather than infrastructure management.

However, asynchronous messaging introduces its own set of challenges. Messages may be delivered multiple times due to the at-least-once delivery model, requiring idempotent consumer design. Visibility timeouts must be carefully configured to balance redelivery responsiveness against the risk of premature message expiry. Dead-letter queues (DLQs) must be tuned to distinguish transient failures from poison messages that will never succeed. Most critically, the configuration parameters that optimize \emph{steady-state performance}—such as batch size for throughput and visibility timeout for latency—may not be optimal when components fail.

\citet{Kyrychenko2025} established a foundational baseline for SQS configuration optimization, demonstrating through queuing theory and empirical measurement that batch size and visibility timeout are the dominant factors affecting throughput and latency under steady load. Their study recommended a batch size of 50, visibility timeout of 600 seconds, and delivery delay of 300 seconds for maximizing throughput in a PDF generation workload processing 500,000 records. These recommendations assume that all components—producers, consumers, and downstream services—function correctly.

\textbf{Real-world systems do not operate in steady state.} Consumer processes crash due to bugs or infrastructure failures. Downstream databases experience transient overload, returning throttling errors or timeouts. Unhandled exceptions occur when messages contain unexpected data or when code defects are triggered by specific inputs. Network partitions isolate components temporarily. In these failure scenarios, the behavior of SQS configuration parameters—particularly visibility timeout and \texttt{maxReceiveCount} (which controls retry attempts before DLQ redrive)—becomes critical. A visibility timeout optimized for low latency in steady state may cause excessive redelivery churn during a downstream outage. A \texttt{maxReceiveCount} set too low may unnecessarily dead-letter messages that would succeed after a brief recovery. Conversely, a \texttt{maxReceiveCount} set too high may waste compute resources on truly broken messages and delay recovery.

Despite the importance of these trade-offs, the academic and practitioner literature provides limited quantitative guidance. AWS documentation offers best-practice recommendations (e.g., "set visibility timeout to expected processing time," "use DLQs to capture failed messages"), but these are qualitative heuristics rather than empirically grounded, statistically validated findings~\citep{AWSResilience2023}. The fault injection literature has focused primarily on infrastructure-level failures such as node crashes and network partitions~\citep{Gao2023CrashFuzz, Wu2024Legolas}, not on application-level faults within queue consumers. The message queue comparison literature evaluates steady-state performance~\citep{DiasDeAssuncao2020MQComparison, Maharjan2023Benchmarking}, not resilience under controlled failure injection.

\section{Research Question}

This research addresses the gap between steady-state optimization and failure resilience by asking:

\begin{quote}
\textbf{How does Amazon SQS configuration affect message reliability and recovery under injected consumer and downstream failures?}
\end{quote}

Specifically, the study investigates:
\begin{itemize}
    \item How visibility timeout, \texttt{maxReceiveCount}, and batch size influence \textbf{message loss}, \textbf{duplicate processing}, and \textbf{dead-letter queue traffic} when consumers fail or downstream services reject requests.
    \item How these configuration parameters affect \textbf{recovery time}—the interval from failure cessation until queue backlog returns to steady state.
    \item Whether the configuration optima identified by \citet{Kyrychenko2025} for steady-state performance transfer to failure scenarios, or whether different trade-offs emerge.
\end{itemize}

\section{Research Objectives}

The objectives of this research are:

\begin{enumerate}
    \item \textbf{Quantify the impact of SQS configuration on message reliability under failure.} Measure message loss (messages that disappear entirely), duplicate processing (messages delivered more than once), and DLQ capture (messages moved to the dead-letter queue) across a matrix of visibility timeout (30–900 seconds), \texttt{maxReceiveCount} (1–10), and batch size (1–50) under four fault types: consumer kill, unhandled error, datastore reject, and datastore timeout.
    
    \item \textbf{Characterize recovery time as a function of configuration parameters.} Define recovery time as the interval from the moment a fault ceases (consumer resumes, downstream recovers) until the queue backlog returns to its pre-failure steady state. Measure how visibility timeout and \texttt{maxReceiveCount} affect this metric.
    
    \item \textbf{Test whether steady-state-optimal configurations remain optimal under failure.} Evaluate whether the configurations recommended by \citet{Kyrychenko2025} for high throughput (large batch size, long visibility timeout) provide the best recovery time, or whether shorter timeouts and different batch sizes perform better in failure scenarios.
    
    \item \textbf{Provide empirical, statistically validated guidance for practitioners.} Conduct pre-registered hypothesis tests with appropriate corrections for multiple comparisons (Holm-Bonferroni), report effect sizes and confidence intervals, and interpret findings in terms of actionable configuration recommendations.
\end{enumerate}

\section{Research Hypotheses}

Three hypotheses are pre-registered and tested:

\begin{description}
    \item[H1 (Message Loss):] Under injected consumer or downstream failures with DLQ configured and sufficient visibility timeout (≥ 900 seconds) and \texttt{maxReceiveCount} (≥ 3), message loss in the SQS queue arm is negligible (< 0.1\%) compared to a synchronous control arm. \emph{Rationale:} SQS redelivery and DLQ capture should prevent loss if configuration allows retries and recovery time.
    
    \item[H2 (maxReceiveCount and Recovery):] After a transient consumer failure (e.g., consumer kill for 10 seconds), recovery time does not differ significantly between \texttt{maxReceiveCount} = 1, 3, 5, and 10 when visibility timeout is held constant at 30 seconds. \emph{Rationale:} Transient faults resolve before retry limits are reached, so \texttt{maxReceiveCount} should not affect recovery duration, only DLQ traffic.
    
    \item[H3 (Steady-State Optimum Transfer):] The configuration that maximizes throughput in steady state (batch 50, visibility timeout 600) has significantly longer recovery time than alternative configurations (batch 10, visibility timeout 30) under consumer failure. \emph{Rationale:} Long visibility timeout delays redelivery; large batches may amplify backlog during failure.
\end{description}

\section{Methodology Overview}

This research employs a controlled experimental design with the following components:

\begin{itemize}
    \item \textbf{Artefact:} An AWS Serverless Application Model (SAM) application implementing a two-arm experiment: a synchronous API arm (control) and an SQS-to-Lambda queue arm (treatment). Both arms process synthetic order payloads, store results in DynamoDB, and emit CloudWatch metrics.
    
    \item \textbf{Fault Injection:} An in-process fault switch integrated into the Lambda consumer handler, capable of injecting four fault types at deterministically configured times: (1) consumer kill (process unavailable), (2) unhandled error (exception thrown), (3) datastore reject (DynamoDB write fails), (4) datastore timeout (delayed write simulating overload).
    
    \item \textbf{Local Simulator:} A discrete-event simulator implementing SQS Standard queue semantics, Lambda event source mapping with batch processing, DynamoDB idempotency, and the fault injection controller on a virtual clock. The simulator enables reproducible, free-tier-aware experimentation without AWS account costs.
    
    \item \textbf{Experiment Runner:} A harness that walks a randomized config matrix, generates manifests with seeds and run IDs, executes pilots and campaigns, aggregates metrics (loss, duplicates, DLQ depth, recovery time, throughput, latency), and persists results in JSON manifests.
    
    \item \textbf{Statistical Analysis:} Hypothesis tests with Holm-Bonferroni correction for family-wise error rate, effect size calculation (rank-biserial r for ordinal comparisons), 95\% bootstrap confidence intervals, and power analysis from pilot runs.
\end{itemize}

All experimental results presented in this submission are generated from the local simulator (`DRY\_RUN=1`), not from live AWS infrastructure. The simulator's discrete-event semantics match SQS behavior (visibility timeout, receive count tracking, DLQ redrive, at-least-once delivery), validated through unit and integration tests. Live AWS validation remains future work; the relative configuration effects measured in simulation are expected to hold, though absolute timing values may differ due to network variability and cold starts.

\section{Contributions}

This research makes the following contributions:

\begin{enumerate}
    \item \textbf{Empirical Evidence on SQS Resilience:} A systematic, quantitative localsim study of how SQS configuration parameters (visibility timeout, \texttt{maxReceiveCount}, batch size) interact with controlled consumer and downstream failures. Prior work evaluated steady-state performance~\citep{Kyrychenko2025} or compared systems without fault injection~\citep{DiasDeAssuncao2020MQComparison}; this work bridges that gap for the present artefact (not a verified global priority claim).
    
    \item \textbf{Reproducible Fault Injection Framework:} A novel local simulator and in-process fault switch enabling deterministic, repeatable injection of application-level faults (consumer kill, unhandled error, datastore reject/timeout) at configurable times. Unlike infrastructure-focused tools~\citep{Gao2023CrashFuzz, Wu2024Legolas}, this approach targets message queue consumer failures specifically.
    
    \item \textbf{Quantitative Configuration Guidance:} Statistical hypothesis tests and effect size measurements provide actionable recommendations for practitioners: when to use short vs. long visibility timeout, how to set \texttt{maxReceiveCount} for transient vs. deterministic faults, and whether steady-state-optimal configurations remain optimal under failure.
    
    \item \textbf{Open Artefact and Dataset:} The entire artefact (SAM template, simulator, experiment runner, analysis scripts, tests, manifests) is committed to the repository with reproducible build instructions. The 350 unique design-cell manifests (690 on-disk files include packaging twins; analysis deduplicates) and 11 figures constitute a reusable dataset for validating alternative analysis approaches or extending to other fault types.
\end{enumerate}

\section{Scope and Limitations}

This study focuses on:
\begin{itemize}
    \item \textbf{SQS Standard queues} (not FIFO, which has different semantics and lower throughput).
    \item \textbf{Lambda consumers} triggered by event source mappings with batch processing.
    \item \textbf{Transient and deterministic-like faults} injected at the consumer/downstream level, not infrastructure failures (AZ outages, SQS service events).
    \item \textbf{Simulated experiments}, not live AWS measurements (though the live backend code is written and tested).
\end{itemize}

Known limitations include:
\begin{itemize}
    \item Cold start latency is not modeled; the simulator assumes warm Lambda execution.
    \item Network variability and jitter are abstracted; the simulator uses deterministic timing.
    \item AWS concurrency limits (soft limit 1,000 concurrent Lambda executions) are not enforced in the simulator.
    \item The study does not evaluate FIFO queues, multi-region architectures, or alternative brokers (Kafka, RabbitMQ).
\end{itemize}

These limitations are documented in \texttt{STATUS.md} and discussed in the conclusion with recommendations for future work.

\section{Document Structure}

The remainder of this document is organized as follows:

\begin{description}
    \item[Chapter 2: Related Work] Surveys the literature on message queue reliability, fault injection, serverless fault tolerance, idempotency, DLQs, recovery metrics, and event-driven architecture. Positions this research within the broader context and identifies the gap it addresses.
    
    \item[Chapter 3: Research Methodology] Describes the experimental design, independent and dependent variables, hypotheses, statistical analysis plan, threats to validity, and the rationale for using a simulator.
    
    \item[Chapter 4: Design and Implementation] Details the SAM architecture, fault injection switch, simulator design, idempotency layer, metrics pipeline, experiment runner, and configuration matrix.
    
    \item[Chapter 5: Evaluation] Presents the experimental results: pilot run repeatability, baseline replication comparison with \citet{Kyrychenko2025}, arms comparison (sync vs. queue), fault campaign results, hypothesis test outcomes with statistical corrections, and figures.
    
    \item[Chapter 6: Conclusions and Discussion] Interprets findings in the context of the research question and objectives, discusses implications for practitioners, evaluates validity, and proposes future work including live AWS validation, FIFO queues, and multi-region scenarios.
\end{description}

The configuration manual, architecture diagrams, demo walkthrough, and viva Q\&A seeds are provided in supplementary documentation (\texttt{docs/}) in the repository.
